\documentclass[12pt]{article}
\usepackage[utf8]{inputenc}
\usepackage[T1]{fontenc}
\usepackage{hyperref}
\title{Las Delicias Zapotec Language Resources for K-12 Curriculum Development}
\author{Compiled by Claude Sonnet 5.5 (claude-sonnet-5-5) for Andr\'as Kornai}
\date{}
\begin{document}
\maketitle

\section{Introduction}
Las Delicias Zapotec (ISO 639-3: zcd) is a Zapotec language of southern Oaxaca, Mexico, in the Otomanguean family. Search results give a population of roughly 16,700 to 17,000 and the vernacular name Didza Xhidza. Glottolog treats it as separate from Rinc\'on Zapotec (zar), although a 2011 source judged some localities easily intelligible with Rinc\'on. The search found only catalogue entries and religious or evangelistic audio materials. No school or linguistic materials were found.

\section{Primary Educational Resources}
No primers, textbooks or school programmes were found in search results.

\section{Linguistic and Reference Materials}
\begin{itemize}
\item Glottolog languoid page for zcd: \url{https://glottolog.org/resource/languoid/iso/zcd}
\item Localizely language-code page for zcd: \url{https://localizely.com/language-code/zcd}
\end{itemize}

\section{Religious and Community Resources}
\begin{itemize}
\item ScriptureEarth listings for zcd: \url{https://www.scriptureearth.org/00eng.php?iso=zcd} and \url{https://www.scriptureearth.org/00i-Scripture_Index.php?iso=zcd}
\item Joshua Project people-group profile (Bible portions and audio Bible reported; no complete Bible or New Testament): \url{https://joshuaproject.net/people_groups/22336}
\item Joshua Project language page: \url{https://joshuaproject.net/index.php/languages/zcd}
\end{itemize}

\section{Digital and Audio Resources}
\begin{itemize}
\item Global Recordings Network page for Zapoteco de Las Delicias: \url{https://globalrecordings.net/en/language/zcd}
\item ``God Came to Us'' recording, MegaVoice: \url{https://megavoice.com/?p=128400}
\end{itemize}

\section{Curriculum Development Context}
Published resources for zcd are scarce. They consist of audio Bible stories and scripture portions. No dictionary, grammar, orthography guide or school material was found. A curriculum would depend on collaboration with speaker communities and on new documentation.
\end{document}